\subsection{Overview}

\textit{[This section will be completed in future work. Placeholder content provided below.]}

The rolling horizon RL architecture consists of three main components:

\begin{enumerate}
    \item \textbf{Policy Network:} Neural network architecture for mapping states to actions
    \item \textbf{Training Algorithm:} PPO, SAC, or Transformer-based RL
    \item \textbf{Rolling Horizon Controller:} Mechanism for sequential re-planning
\end{enumerate}

\subsection{Policy Network Architecture}

\textit{[To be completed: Network architecture details, including:]}
\begin{itemize}
    \item Input layer: State representation (project features, portfolio status)
    \item Hidden layers: MLP, LSTM, or Transformer encoder
    \item Output layer: Budget allocation distribution
    \item Attention mechanism for project prioritization (optional)
\end{itemize}

\subsection{Training Algorithm}

\textit{[To be completed: Algorithm selection and hyperparameters, including:]}
\begin{itemize}
    \item Candidate algorithms: PPO, SAC, A3C
    \item Hyperparameter tuning methodology
    \item Training dataset size and diversity
    \item Convergence criteria
\end{itemize}

\subsection{Rolling Horizon Implementation}

\textit{[To be completed: Implementation details, including:]}
\begin{itemize}
    \item Episode generation for fixed $H=12$ horizon
    \item State initialization and masking for inactive projects
    \item Action feasibility enforcement
    \item Reward shaping and normalization
\end{itemize}

\subsection{Computational Considerations}

\textit{[To be completed: Computational analysis, including:]}
\begin{itemize}
    \item Training time and resource requirements
    \item Inference latency (target: <200ms per decision)
    \item Scalability to larger portfolios ($N > 20$)
    \item GPU vs. CPU performance comparison
\end{itemize}
